\section{Resultados}

\subsection{Análisis Exploratorio}
El análisis de correlación reveló relaciones significativas entre las
características de audio: energy y loudness presentan correlación positiva
fuerte ($r > 0.7$), mientras que energy y acousticness muestran correlación
negativa ($r < -0.7$), consistente con la naturaleza opuesta de música
electrónica/acústica \cite{schedl2018current}.

\subsection{Determinación del K Óptimo}
% INSTRUCCIÓN: Completar con resultados reales del notebook ejecutado
El método del codo sugirió $K = $ [valor] como punto de inflexión.
El Silhouette Score máximo se obtuvo en $K = $ [valor] con un valor de [valor].

\subsection{Caracterización de Clusters}
% INSTRUCCIÓN: Completar tabla con valores reales del perfil promedio
\begin{table}[H]
\centering
\caption{Perfil Promedio de Características de Audio por Cluster}
\label{tab:clusters}
\begin{tabular}{lccccc}
\toprule
\textbf{Cluster} & \textbf{Dance.} & \textbf{Energy} & \textbf{Valence} & \textbf{Acoust.} & \textbf{N} \\
\midrule
0 & -- & -- & -- & -- & -- \\
1 & -- & -- & -- & -- & -- \\
2 & -- & -- & -- & -- & -- \\
3 & -- & -- & -- & -- & -- \\
4 & -- & -- & -- & -- & -- \\
5 & -- & -- & -- & -- & -- \\
\bottomrule
\end{tabular}
\end{table}

\subsection{Visualización PCA}
La proyección de los clusters en el espacio de 2 componentes principales
(ver figura \ref{fig:pca}) permitió observar la separación entre grupos.
La varianza explicada acumulada por las dos primeras componentes fue de
[X]\%.
\begin{figure}[H]
\centering
\includegraphics[width=\columnwidth]{figuras/pca_clusters.png}
\caption{Visualización de clusters K-Means mediante PCA 2D}
\label{fig:pca}
\end{figure}
